\subsection{Context width: compiling any circuit into a DNNF}\label{sec:ctxwidth}

DFSB branches on shared variables but treats overlapping $\vee$'s by envelopes.
With the \#DNNF FPRAS available, non-deterministic $\vee$'s cost nothing, and
the only obstacle is $\wedge$-gates whose children share variables. The following
parameter measures exactly that, and generalises Darwiche's dtree conditioning
\cite{darwiche2001dnnf} from CNFs to arbitrary circuits: a binary $\wedge$-tree over
the clauses of a CNF is a dtree, and $\kappa$ of its internal nodes below is their
cluster (cutset plus context). Like the width of a dtree, it is a parameter of how
the $\wedge$'s are bracketed, not of the function.

\begin{definition}[context width]
Let $N$ be the NNF of $C$ (\cref{sec:setting}) in which every $\wedge$-gate keeps the
children it has in $C$: nested $\wedge$-gates are not merged, and the children of a
gate form a set ($h\wedge h=h$). $|N|$ is its number of wires. For an $\wedge$-gate
$g$ let $S_g=\bigcup_{a\ne b\text{ children of }g}\sv(a)\cap\sv(b)$ be its overlap.
For a gate $h$ let $\mathrm{Up}(h)$ be the set of $\wedge$-gates from which $h$ is
reachable (including $h$), and define the \emph{context}
$\kappa(h)=\sv(h)\cap\bigcup_{g\in\mathrm{Up}(h)}S_g$. The context width is
$\mathrm{cw}(N)=\max_h|\kappa(h)|$.
\end{definition}

For $\{\wedge,\vee,\neg\}$-circuits $|N|\le2|C|$ (one copy of each gate per
polarity). A fan-in-$k$ $\oplus$-gate becomes a chain with $12(k-1)$ wires for both
polarities, so $|N|\le12|C|$ when $\oplus$-gates are present, and a threshold gate of
fan-in $k$ with nonnegative integer weights of total $W$ becomes a sequential
counter with $O(kW)$ wires; other threshold gates are covered only as frontier
gates (see below). Merging an $\vee$-gate into an $\vee$-parent does not increase
$\mathrm{cw}$ ($\vee$-gates have no overlap, and the sets $\mathrm{Up}(h)$ do not
change), but merging $\wedge$-gates can: a CNF with clauses on the $n-w+1$ sliding
windows of width $w$ of its variable order has $\mathrm{cw}=w$ as a left-deep
$\wedge$-chain in window order, and $\mathrm{cw}=n-2$ (for $2\le w\le n-2$) as a single
$\wedge$-gate over its clauses. We write $\mathrm{cw}^*(C)$ for the minimum of
$\mathrm{cw}$ over the NNFs obtained from $N$ by merging $\wedge$-gates of fan-out $1$
into their $\wedge$-parents and replacing $\wedge$-gates by binary trees of
$\wedge$-gates over the same children (which at most doubles $|N|$); any bracketing an
algorithm chooses, e.g.\ that of \cref{lem:pathcw}, bounds it.

\begin{theorem}[context-width compilation]\label{thm:cw}
Let $N$ be an NNF of $C$ as above and $\mathrm{cw}=\mathrm{cw}(N)$. Then $C$ is
equivalent to a DNNF with at most $(|N|+1)(\mathrm{cw}+3)2^{\mathrm{cw}}$ wires,
computable in time $O(|N|(n+(\mathrm{cw}+1)2^{\mathrm{cw}}))$. $\mathrm{cw}(N)=0$ iff
$N$ is a DNNF. If every $\vee$-gate of $N$ is deterministic (those produced by the
$\oplus$- and threshold expansions are; those produced from $\wedge$-gates of $C$ under
a negation need not be), the result is a d-DNNF, and $p$ is computed exactly in time
linear in its size under every product distribution. Consequently (via
\cref{thm:sis}(a) and \cite{meel2026dnnf}) $p$ has an FPRAS under every product
distribution with rational marginals $q_i=a_i/b_i$, running in time
$\poly(|N|2^{\mathrm{cw}},n,\max_i\log b_i,1/\eps,\log(1/\delta))$; polynomial when
$\mathrm{cw}=O(\log|N|)$.
\end{theorem}

\begin{proof}
Two facts. (F1) If $h$ is a child of $g$ then $\mathrm{Up}(h)\supseteq\mathrm{Up}(g)$,
so $\kappa(g)\cap\sv(h)\subseteq\kappa(h)$. (F2) If $h_i,h_j$ are distinct children
of an $\wedge$-gate $g$ then $\sv(h_i)\cap\sv(h_j)\subseteq S_g\cap\sv(g)\subseteq\kappa(g)$.
Nodes of the DNNF: for each gate $h$ and $\sigma\in\bits^{\kappa(h)}$ a node
$N[h,\sigma]$ for $h|_\sigma$; for each wire $g\to h$ and each assignment $\pi$ to
$\kappa(g)\cap\sv(h)$ a node
$M[h,\pi]=\bigvee_{\tau}\big(\mathrm{lits}(\tau)\wedge N[h,\pi\cup\tau]\big)$, $\tau$
ranging over assignments to $\kappa(h)\setminus\mathrm{dom}\,\pi$ (a deterministic
Shannon expansion; $\pi\cup\tau$ covers $\kappa(h)$ by F1). Set
$N[g,\sigma]=\mathrm{op}_j M[h_j,\sigma|_{\kappa(g)\cap\sv(h_j)}]$ for a gate $g$ with
children $h_j$, literal leaves become constants when their variable is in
$\kappa(h)$, and the root is $M[\mathrm{root},\emptyset]$. Correctness is induction
on $g|_\sigma=\mathrm{op}_jh_j|_{\sigma}$. Decomposability: inside $M$, the literals
of $\tau$ live on $\kappa(h)\setminus\mathrm{dom}\,\pi$ and $N[h,\cdot]$ on
$\sv(h)\setminus\kappa(h)$; at an $\wedge$-gate the children's variables lie in
$\sv(h_j)\setminus\kappa(g)$, pairwise disjoint by F2.
Size: a gate $h$ gives $2^{|\kappa(h)|}$ nodes $N[h,\cdot]$ with $\mathrm{fanin}(h)$
wires each, at most $|N|2^{\mathrm{cw}}$ wires in all. A wire $g\to h$ gives
$2^{|\mathrm{dom}\,\pi|}$ nodes $M[h,\pi]$ with at most $(t+2)2^t$ wires each,
$t=|\kappa(h)|-|\mathrm{dom}\,\pi|$, so at most $(\mathrm{cw}+2)2^{\mathrm{cw}}$; the root
adds as much. Time: supports, overlaps and contexts take $O(n)$ bit operations per
wire of $N$, and each wire of the output is built in time $O(\mathrm{cw}+1)$.
Since $S_g\subseteq\kappa(g)$, $\mathrm{cw}=0$ iff every $S_g$ is empty.
The d-DNNF claim holds because conditioning preserves disjointness and the
$M$-nodes are deterministic.
\end{proof}

\begin{lemma}[bracketing along a path decomposition]\label{lem:pathcw}
Let $F$ be a CNF whose primal graph (variables adjacent when they share a clause)
has a path decomposition $B_1,\dots,B_r$ of width $k$. Let $\ell(c)$ be the index of
the first bag containing $\sv(c)$ (one exists since $\sv(c)$ is a clique), order the
clauses as $c_1,\dots,c_m$ with $\ell(c_1)\le\dots\le\ell(c_m)$, and bracket $F$ as the
left-deep chain $a_1=c_1$, $a_j=a_{j-1}\wedge c_j$. Then $\kappa(a_j)$ and
$\kappa(c_j)$ lie in $B_{\ell(c_j)}$, so $\mathrm{cw}\le k+1$. The same bound holds for
an $\vee$ of such CNFs on the same variables whose chains share no $\wedge$-gate. For
constant $k$, a path decomposition of width $k$, if one exists, is found in linear
time \cite{bodlaender1996,bodlaenderkloks1996}, and the order from it in linear time.
\end{lemma}

\begin{proof}
Write $\ell_j=\ell(c_j)$. The bags containing a variable form an interval, so if
$v\in\sv(c_i)\cap\sv(c_{i'})$ with $i\le j\le i'$ then $v\in B_{\ell_i}\cap B_{\ell_{i'}}
\subseteq B_{\ell_j}$. The $\wedge$-gates above $a_j$ are $a_{j+1},\dots,a_m$, and
$S_{a_{j'}}\subseteq\sv(c_{j'})$. So $v\in\kappa(a_j)$ means $v\in\sv(c_i)$ for some
$i\le j$ and $v\in\sv(c_{j'})$ for some $j'\ge j$, hence $v\in B_{\ell_j}$. Further
$\kappa(c_j)\subseteq\sv(c_j)\subseteq B_{\ell_j}$, a literal has $|\kappa|\le1$, and
bags have at most $k+1$ elements. In an $\vee$ of such CNFs, $\vee$-gates add no
overlap and the $\wedge$-gates above $a_j$ are those of its own chain; a clause or
literal shared between chains still has $\kappa$ inside its own variables. For the
running time: clauses have at most $k+1$ variables, so the primal graph has
$O(k|F|)$ edges, and $\ell(c)$ is the largest index of a first bag of a variable
of $c$.
\end{proof}

\begin{example}[$\vee$'s are easy, $\wedge$'s are hard]\label{ex:or-cnf}
Let $C=\bigvee_{i\le t}F_i$ where each $F_i$ is a CNF with clauses on sliding
windows of width $w$ along its own variable order, bracketed as a left-deep
$\wedge$-chain in window order. The windows are a path decomposition of width $w-1$
of each $F_i$, so $\mathrm{cw}\le w$ by \cref{lem:pathcw}, and an FPRAS exists when
$w=O(\log n)$, while $p$ can be $2^{-\Theta(n)}$. Exact computation is $\#\mathsf
P$-hard already for $\mathrm{cw}=0$, for ORs of at most $5$ \emph{read-once} CNFs:
greedily edge-colour a graph of maximum degree $3$ with $5$ colours, write its
vertex-cover CNF as $\bigwedge_{c\le5}\varphi_c$ with each $\varphi_c$ a matching
(hence read-once) CNF, and use
$\#\bigwedge_c\varphi_c=\sum_{\emptyset\ne S\subseteq[5]}(-1)^{|S|+1}\#\bigvee_{c\in S}\varphi_c$;
counting vertex covers (equivalently, independent sets) is $\#\mathsf P$-complete
in graphs of maximum degree $3$ \cite{greenhill2000}.
\end{example}

The same construction works with a \emph{frontier} of gates that are compiled
otherwise: gates of small support by tabulation (as the implementation does), and
circuits of treewidth $k$ into d-SDNNFs of size $2^{O(k)}|C|$ \cite{acms2020}.
Context width is then computed with frontier gates (one per polarity) as leaves; a
frontier gate $f$ compiled to a DNNF $D_f$ gives the nodes $N[f,\sigma]=D_f|_\sigma$,
so the size becomes $\sum_f2^{|\kappa(f)|}|D_f|$ plus the bound of \cref{thm:cw} for
the remaining gates, and the result is a d-DNNF only if every $D_f$ is one.

\subsection{Modules and quotients}\label{sec:modules-quotients}

A gate $g$ in the cone of the output is a \emph{module} if every path from an input
in $\sv(g)$ to the output passes through $g$. Then $C(x)=Q_g(g(x_{\sv(g)}),x_{\rm rest})$
for the \emph{quotient} $Q_g$ obtained by replacing $g$ with a fresh input, $x_{\rm
rest}$ being the inputs outside $\sv(g)$, and under a product distribution $g(x)$ is
independent of $x_{\rm rest}$. The supports of modules form a laminar family: if
$v\in\sv(g)\cap\sv(h)$, a path from $x_v$ through $g$ to the output passes through
$h$, so one of $g,h$ is in the cone of the other and their supports are nested.
Modules are found in linear time \cite{dutuitrauzy1996}. The maximal modules
$g_1,\dots,g_k$ strictly inside a module $g$ have disjoint supports, and $g$ is a
function $Q_g(g_1,\dots,g_k,x_{\rm rest})$ of them and of its remaining inputs; this
is the module tree.

\begin{lemma}[sensitivity of quotients]\label{lem:quotient}
Let $F(q)=\Pr_{y\sim\bigotimes_j\mathrm{Bern}(q_j)}[Q(y)=1]$ and $q,\hat q\in[0,1]^k$.
If $\hat q_j\in e^{\pm\xi_j}q_j$ and $1-\hat q_j\in e^{\pm\xi_j}(1-q_j)$ for all $j$,
then $F(\hat q)\in e^{\pm\sum_j\xi_j}F(q)$. If $Q$ is monotone (resp.\ antitone) in
$y_j$, the condition on $q_j$ (resp.\ $1-q_j$) alone suffices for that coordinate.
If $q_j\in[L_j,U_j]$ with $U_j\le\frac12$, then $\hat q_j=\sqrt{L_jU_j}$ satisfies both
conditions with $\xi_j=\frac12\log(U_j/L_j)$ (and symmetrically, for $1-q_j$, if
$L_j\ge\frac12$).
\end{lemma}

\begin{proof}
$F$ is multilinear: $F=q_jA+(1-q_j)B$, with $A,B\ge0$ when the other coordinates
lie in $[0,1]$, so scaling either term by a factor in $e^{\pm\xi_j}$ scales $F$ within
$e^{\pm\xi_j}$. If $Q$ is monotone in $y_j$ then $A\ge B$, and for $c\ge1$ with
$cq_j\le1$, $F\le F(cq_j)=F+(c-1)q_j(A-B)\le F+(c-1)q_jA\le cF$; symmetrically for
$c\le1$. Change one coordinate at a time; every intermediate point lies in
$[0,1]^k$. For the last claim let $r=\sqrt{U_j/L_j}$, so $\hat q_j=rL_j=U_j/r$ and
$\hat q_j/q_j\in[1/r,r]$. The ratio $(1-\hat q_j)/(1-q_j)$ lies between
$(1-\hat q_j)/(1-L_j)$ and $(1-\hat q_j)/(1-U_j)$, both in $[1/r,r]$ because
$r(1-\hat q_j)-(1-L_j)=(r-1)(1-L_j-\hat q_j)\ge0$ and
$r(1-U_j)-(1-\hat q_j)=(r-1)(1-U_j-\hat q_j)\ge0$ when $U_j\le\frac12$.
\end{proof}

Let $g$ be a module with maximal submodules $g_1,\dots,g_k$, and let
$q_j=D(g_j)\in[L_j,U_j]$. Take $\hat q_j,\xi_j$ from the last part of
\cref{lem:quotient} when $U_j\le\frac12$ or $L_j\ge\frac12$, and otherwise any
$\hat q_j\in[L_j,U_j]$ and $\xi_j=\max\{\log(U_j/L_j),\log((1-L_j)/(1-U_j))\}$ (only
the first term if $Q_g$ is monotone in $y_j$, only the second if it is antitone).
An interval $[L_F,U_F]$ for $F(\hat q)$, i.e.\ for $Q_g$ under
$\bigotimes_j\mathrm{Bern}(\hat q_j)$ and $D$ on $x_{\rm rest}$, from any sound method
(exact for small quotients, the threshold dynamic program, DFSB, Monte Carlo), gives
$D(g)\in[e^{-\sum_j\xi_j}L_F,e^{\sum_j\xi_j}U_F]$; an interval for $1-F(\hat q)$ (that
is, for $\neg Q_g$) gives one for $D(\neg g)$ in the same way. Walking the module tree
bottom-up and making each of the $M$ quotient estimates at confidence $\delta/M$
gives an interval for $p$ that is valid with probability $1-\delta$ (union bound; each
estimate is valid conditionally on the earlier ones). If at each module the rarer
polarity of the quotient is estimated and every interval lies on one side of
$\frac12$, the log-ratio $\log(U/L)$ of the final interval is the sum of the
log-ratios of all quotient estimates. Completeness thus needs relative accuracy on
the rarer of $D(g)$ and $D(\neg g)$ at every module; if every quotient is monotone or
antitone in each $y_j$, it suffices to estimate at each module the polarity its
parent's condition needs. Thresholds, majorities and other tractable functions of
independent sub-circuits are the typical case.
